\documentclass[11pt]{article}
\usepackage[letterpaper,margin=0.88in]{geometry}
\usepackage{amsmath,amssymb,amsthm,booktabs,graphicx,microtype,xcolor}
\usepackage[colorlinks=true,linkcolor=blue!55!black,citecolor=blue!55!black,urlcolor=blue!55!black]{hyperref}

\newtheorem{theorem}{Theorem}
\newtheorem{lemma}{Lemma}
\newcommand{\R}{\mathbb R}
\newcommand{\smin}{s_{\min}}
\newcommand{\sharpw}{\mathbin{\sharp}}

\title{A counterexample to a weighted-geometric-mean\\log-majorization conjecture}
\author{Candidate counterexample for arXiv:2105.13356}
\date{11 August 2026}

\begin{document}
\maketitle

\noindent\fbox{\parbox{0.95\linewidth}{\textbf{Status.}
Candidate counterexample, likely valid, requiring human review.  Exact integer
positive-definite matrices disprove Conjecture 1.1 of Ghabries--Abbas--Mourad--Assi.
The decisive numerical inequality is certified by exact rational ball arithmetic.}}

\section{Source conjecture}

For positive-definite matrices $A,B$, their weighted geometric mean is
\[
 A\sharpw_t B=A^{1/2}(A^{-1/2}BA^{-1/2})^tA^{1/2},\qquad 0\le t\le1.
\]
Writing $s_1(Y)\ge\cdots\ge s_n(Y)$ for the singular values of $Y$, one says
$s(X)\prec_{\log}s(Y)$ when
\[
 \prod_{j=1}^k s_j(X)\le \prod_{j=1}^k s_j(Y)\quad(1\le k<n),
 \qquad |\det X|=|\det Y|.
\]
Conjecture 1.1 in Ghabries--Abbas--Mourad--Assi~\cite{GAMA} asks whether
\begin{equation}\label{eq:conjecture}
 s\!\left(A^t(A\sharpw_tB)B^{1-t}\right)\prec_{\log}s(AB)
 \qquad(0\le t\le1)
\end{equation}
for all positive semidefinite $A,B$.  It originates in work of Lemos and
Soares~\cite{LS}.  The source statement is reproduced below.

\begin{figure}[ht]
\centering
\includegraphics[width=0.91\linewidth]{figures/open_problem_crop.png}
\caption{Conjecture 1.1 in the source PDF~\cite[p.~3]{GAMA}.}
\end{figure}

\section{Idea of the proof}

The determinant in \eqref{eq:conjecture} always agrees on the two sides.
Consequently, in dimension four its third prefix-product inequality says
exactly that the smallest singular value of the left-hand matrix is at least
the smallest singular value of $AB$.  This dual endpoint is where the
conjecture fails.

The example below makes one matrix diagonal and uses a rational value
$t=15/16$, so all fractional powers can be built from four successive positive
square roots.  Exact $LDL^{\mathsf T}$ arithmetic proves a lower bound for
$\smin(AB)$.  Repeated rational matrix balls enclose the square roots and give
an upper bound for the other smallest singular value.  The two certified
bounds are separated by more than $0.08$.

\section{The counterexample}

\begin{theorem}\label{thm:counterexample}
Let $t=15/16$ and let
\[
A=\begin{pmatrix}
17825066&0&0&0\\
0&83&0&0\\
0&0&24&0\\
0&0&0&10
\end{pmatrix},
\qquad
B=\begin{pmatrix}
117567&-99513&2310&-38624\\
-99513&131686&572&24611\\
2310&572&190&-1191\\
-38624&24611&-1191&14107
\end{pmatrix}.
\]
Then $A$ and $B$ are positive definite, but
\[
s\!\left(A^t(A\sharpw_tB)B^{1-t}\right)
\not\prec_{\log}s(AB).
\]
More precisely, if $X=A^t(A\sharpw_tB)B^{1-t}$, then
\begin{equation}\label{eq:gap}
 \smin(X)<235.715744<235.8<\smin(AB).
\end{equation}
Numerically, the violated third-prefix ratio is
\[
 \frac{\prod_{j=1}^3s_j(X)}{\prod_{j=1}^3s_j(AB)}
 =1.0009342385069400\ldots>1.
\]
\end{theorem}

\begin{proof}
The positivity of $A$ is immediate.  The four leading principal minors of $B$
are
\[
117567,\quad 5579090793,\quad 55894510422,\quad 2308323855975,
\]
so $B>0$ by Sylvester's criterion.

For arbitrary positive-definite $A,B$,
\[
 \det(A\sharpw_tB)=(\det A)^{1-t}(\det B)^t.
\]
It follows that
\begin{equation}\label{eq:det}
 \det X=(\det A)^t\det(A\sharpw_tB)(\det B)^{1-t}
 =\det A\det B=\det(AB).
\end{equation}
Because all singular values are positive, \eqref{eq:det} gives
\[
 \frac{\prod_{j=1}^3s_j(X)}{\prod_{j=1}^3s_j(AB)}
 =\frac{\smin(AB)}{\smin(X)}.
\]
Thus it remains only to prove \eqref{eq:gap}.

Set $q=1179/5=235.8$.  Exact integer multiplication gives that
\[
 D=25(AB)^{\mathsf T}(AB)-1179^2I_4
\]
has leading principal minors
\begin{align*}
\Delta_1&=109792626457520745575173084,\\
\Delta_2&=42609956693745032985833310177535659834256,\\
\Delta_3&=2228978970825693898612895858077918124109611904,\\
\Delta_4&=182721114621736343069686961423087938455547419076911.
\end{align*}
All are positive.  Hence $D>0$ and therefore
$(AB)^{\mathsf T}(AB)>q^2I_4$, proving $\smin(AB)>q$.

We now certify the other half of \eqref{eq:gap}.  The included verifier uses
operator-norm balls with exact rational centers and radii.  Its square-root
step rests on the following elementary bound.  If $S,R>0$, then
\begin{equation}\label{eq:sqrtbound}
 \|R-S^{1/2}\|_F
 \le
 \frac{\|R^2-S\|_F}{\lambda_{\min}(R)+\sqrt{\lambda_{\min}(S)}}.
\end{equation}
Indeed, for $E=R-S^{1/2}$ one has
$RE+ES^{1/2}=R^2-S$; the Sylvester operator on the left has least singular
value at least
$\lambda_{\min}(R)+\sqrt{\lambda_{\min}(S)}$ in the Frobenius norm.

Starting from the exact matrices $A,B$, four applications of
\eqref{eq:sqrtbound} enclose the $1/16$ powers of $A$, $B$, and
$A^{-1/2}BA^{-1/2}$.  Products and inverses use the standard norm-ball rules
\[
 e_{UV}\le\|U_0\|e_V+\|V_0\|e_U+e_Ue_V,
 \qquad
 e_{U^{-1}}\le
 \frac{\|U_0^{-1}\|^2e_U}{1-\|U_0^{-1}\|e_U}.
\]
All displayed norms are bounded above by exact rational Frobenius bounds, and
all eigenvalue lower bounds use exact rational inverses and $LDL^{\mathsf T}$
pivots.  The resulting ball for $X$ has radius less than
$2.546\times10^{-42}$.  An exact rational trial vector $v$ then gives
\[
 \smin(X)\le \frac{\|Xv\|}{\|v\|}<235.715744.
\]
Every comparison in this last paragraph is executed as an exact rational
inequality by \path{code/certify_counterexample.py}; floating point is used
only to select the rational centers and trial vector.  This proves
\eqref{eq:gap}, hence the third prefix inequality is reversed and
\eqref{eq:conjecture} is false.
\end{proof}

\section{Verification and scope}

Running
\begin{quote}\ttfamily\small
conda run --no-capture-output -n sandbox python
code/certify\_counterexample.py
\end{quote}
reproduces the rational certificate.  Its high-precision centers give
\[
\smin(X)=235.715743188457775258697\ldots,
\quad
\smin(AB)=235.935957912436422532775\ldots.
\]
The separate search script records the heuristic discovery route; it is not
used in the proof.

The counterexample is strict, so continuity yields an open interval of failing
parameters around $15/16$.  No claim is made that dimension four is minimal,
and the source's other conjectures are not decided here.

\paragraph{Bounded novelty check.}
The run's registry, solution, attempt, and proof-gap indexes; exact expressions;
the source title, authors, and DOI; weighted geometric-mean and
log-majorization keywords; and the source's citing works were searched through
11 August 2026.  Gan and Kim reproduce the exact conjecture and call it still
open~\cite[Sec.~3]{GK}.  The three indexed citing papers through that date do
not state a decision of this conjecture, and no prior counterexample was found.
This is a bounded audit, not a definitive priority claim.

\paragraph{Human-review recommendation.}
First audit the exact matrix transcription and the determinant reduction
\eqref{eq:det}.  Then inspect the square-root enclosure
\eqref{eq:sqrtbound}, the propagation of ball radii, and the final exact
rational assertions in the verifier.  These are the only computer-assisted
parts of the proof.

\begin{thebibliography}{9}
\bibitem{GAMA}
M.~M. Ghabries, H.~Abbas, B.~Mourad, and A.~Assi,
\emph{New log-majorization results concerning eigenvalues and singular values
and a complement of a norm inequality}, Linear Multilinear Algebra 71 (2023),
1228--1243; arXiv:2105.13356.

\bibitem{LS}
R.~Lemos and G.~Soares,
\emph{Some log-majorizations and an extension of a determinantal inequality},
Linear Algebra Appl. 547 (2018), 19--31.

\bibitem{GK}
L.~Gan and S.~Kim,
\emph{Weak log-majorization between the geometric and Wasserstein means},
J. Math. Anal. Appl. 530 (2024), 127711; arXiv:2301.07934.
\end{thebibliography}

\end{document}
